\documentclass[9pt,a4paper,landscape]{article}
\usepackage[top=1.4cm,bottom=1.4cm,left=1.8cm,right=1.8cm]{geometry}
\usepackage{fontenc}
\usepackage{palatino}
\usepackage{booktabs}
\usepackage{xcolor}
\usepackage{tabularx}
\usepackage{enumitem}
\usepackage{parskip}
\usepackage{array}
\usepackage{colortbl}

% Stage colours
\definecolor{sparse}{RGB}{120,120,120}
\definecolor{building}{RGB}{0,160,180}
\definecolor{peak}{RGB}{200,50,40}
\definecolor{recap}{RGB}{50,160,80}
\definecolor{resolution}{RGB}{200,160,0}

\definecolor{sparsebg}{RGB}{235,235,235}
\definecolor{buildingbg}{RGB}{220,245,248}
\definecolor{peakbg}{RGB}{252,228,227}
\definecolor{recapbg}{RGB}{224,245,228}
\definecolor{resolutionbg}{RGB}{252,245,220}

\setlist[itemize]{leftmargin=1.2em, itemsep=0pt, topsep=1pt, parsep=0pt}

\setlength{\parskip}{0pt}
\setlength{\parindent}{0pt}

\newcommand{\stageheader}[3]{%
  \cellcolor{#1}\textcolor{white}{\textbf{#2}}\\
  \cellcolor{#1}\textcolor{white}{\small #3}%
}

\begin{document}

\begin{center}
  {\large\bfseries Wolfson — Performance Plan}\quad
  {\small Bass improvisation cues by arc stage \;|\; 1 May 2026, Wolfson College Oxford}
\end{center}

\vspace{3pt}

\begin{tabularx}{\linewidth}{@{}|X|X|X|X|X|@{}}

% Header row
\cellcolor{sparse}\textcolor{white}{\textbf{Sparse}}~\newline\cellcolor{sparse}\textcolor{white}{\small 0:00–1:00} &
\cellcolor{building}\textcolor{white}{\textbf{Building}}~\newline\cellcolor{building}\textcolor{white}{\small 1:00–2:30} &
\cellcolor{peak}\textcolor{white}{\textbf{Peak}}~\newline\cellcolor{peak}\textcolor{white}{\small 2:30–3:30} &
\cellcolor{recap}\textcolor{white}{\textbf{Recapitulation}}~\newline\cellcolor{recap}\textcolor{white}{\small 3:30–4:30} &
\cellcolor{resolution}\textcolor{white}{\textbf{Resolution}}~\newline\cellcolor{resolution}\textcolor{white}{\small 4:30–5:00} \\

\hline

% Content row
\cellcolor{sparsebg}
\begin{itemize}
  \item \textbf{Phrases echo back literally} — same intervals and rhythm, transposed to sax register. Play short, distinctive phrases.
  \item \textbf{Leave 2+ beats of rest} after each phrase — system needs \textasciitilde2\,s silence to close the phrase and respond.
  \item \textbf{Play quietly} — sax mirrors your dynamic. A quiet opening makes later contrast dramatic.
  \item \textbf{Plant a 2–3 note motif} — it will be recalled in recapitulation.
\end{itemize}
&
\cellcolor{buildingbg}
\begin{itemize}
  \item \textbf{Play D Dorian} (D E F G A B C) — 4+ distinct pitch classes shifts the sax harmonic language. Switch away briefly to hear it change.
  \item \textbf{Repeat a riff 3+ times} — sax enters development mode: stronger motivic use, directed contour. It starts to \textit{argue with} the riff.
  \item \textbf{Play continuously (\textasciitilde8\,s)} to trigger trading mode — sax breaks in, then re-enters every \textasciitilde4\,s (just under a bar).
  \item \textbf{Let the sax repeat} (\texttt{-\,-sax-riff-prob}) — sax insists on its phrase; you develop underneath. After two replays it varies itself.
\end{itemize}
&
\cellcolor{peakbg}
\begin{itemize}
  \item \textbf{Play loudly and densely, then sparsely} — busy bass invites sparser sax; a terse phrase invites a rush of notes.
  \item \textbf{Rising phrase then falling} — contour steering. Sax responds to a rising line with a descending resolution and vice versa.
  \item \textbf{Switch between Q\,\&\,A and continuous playing} — short phrases with 2-beat rests give call-and-response; sustained riff triggers trading mode. The contrast is most audible here.
\end{itemize}
&
\cellcolor{recapbg}
\begin{itemize}
  \item \textbf{Play low} — register contrast pushes sax into upper register. Spatial separation is most audible here.
  \item \textbf{Listen} — sax quotes motifs from sparse and building stages. Pause to let the audience hear the system remembering.
\end{itemize}
&
\cellcolor{resolutionbg}
\begin{itemize}
  \item \textbf{Stop playing} — sax closes the piece unaccompanied. Put the bass down deliberately. The clearest demonstration of musical agency.
\end{itemize}
\\

\hline
\end{tabularx}

\vspace{4pt}
{\small
\textbf{Q\,\&\,A mode:} play a phrase, stop for 2+ beats — sax replies \textasciitilde1\,s after silence.\quad
\textbf{Trading mode:} play continuously for \textasciitilde8\,s — sax breaks in and re-enters every \textasciitilde4\,s.\quad
\textbf{Echo:} first 4 exchanges are literal echoes; after that Wolfson generates freely.\quad
\textbf{Motif:} repeat sparse motif 2--3$\times$ for strong recapitulation recall.\quad
\textbf{Auto-start:} Wolfson plays first; join when ready.
}

\end{document}
